\input{configTD}

% ── Poly à trous : commandes spécifiques ──────────────────────────

\newcommand{\atrou}[2][3cm]{%
  \ifthenelse{\boolean{showSolutions}}{%
    \par\vspace{0.5em}%
    \begin{mdframed}[backgroundcolor=ThemeLight, linewidth=0pt,
      skipabove=0pt, skipbelow=0pt,
      innertopmargin=8pt, innerbottommargin=8pt]%
    #2%
    \end{mdframed}%
    \par\vspace{0.5em}%
  }{%
    \par\vspace{0.5em}%
    \noindent\fbox{\begin{minipage}[c][#1]%
      {\dimexpr\textwidth-2\fboxsep-2\fboxrule\relax}%
      \color{white}\rule{0pt}{0pt}%
    \end{minipage}}%
    \par\vspace{0.5em}%
  }%
}

\newcommand{\val}[1]{%
  \ifthenelse{\boolean{showSolutions}}{#1}{\phantom{#1}}%
}

\newcommand{\esquisse}[1][5.5cm]{%
  \begin{center}%
  \ifthenelse{\boolean{showSolutions}}{%
    \fbox{\begin{minipage}[c][#1]{0.85\textwidth}%
      \centering\vfill
      \textit{\small Croquis réalisé en cours.}%
      \vfill
    \end{minipage}}%
  }{%
    \fbox{\begin{minipage}[c][#1]{0.85\textwidth}%
      \hfill\vfill\hfill
    \end{minipage}}%
  }%
  \end{center}%
}

\newcommand{\degC}{\text{\textdegree C}}

\title{\sffamily\bfseries Analyse Numérique --- CM2\\[0.3em]
  \Large Différences finies, systèmes linéaires,\\
  stabilité et convergence}
\author{}
\date{19 mars 2026}

\begin{document}
\sffamily
\maketitle
\thispagestyle{empty}

% ====================================================================
\section*{1. Rappel --- Ce qu'on a retenu du CM1}
% ====================================================================

Au CM1, on a construit une méthode de calcul approché
(méthode d'Euler) à partir du refroidissement d'un café.

\medskip

\textbf{Questions de rappel.}
\begin{itemize}[leftmargin=*]
  \item Qu'a-t-on remplacé dans la méthode d'Euler ?
  \item Pourquoi faut-il choisir un pas ?
  \item Qu'est-ce qui peut rendre une approximation peu fiable ?
\end{itemize}

\atrou[3.6cm]{%
\begin{itemize}[leftmargin=*]
  \item On a remplacé la \textbf{dérivée} (information continue)
        par un \textbf{quotient de différences} (information discrète).
  \item Le pas $\Delta t$ contrôle la \textbf{précision} :
        plus il est petit, plus on suit fidèlement la courbe exacte,
        mais plus le calcul coûte cher.
  \item Un pas trop grand peut produire des résultats
        \textbf{physiquement absurdes}
        (température négative, divergence).
\end{itemize}
}

\medskip

\textbf{Question de transition.}

\begin{center}
\begin{mdframed}[backgroundcolor=ThemeLight, linecolor=Theme,
  linewidth=1pt, innertopmargin=8pt, innerbottommargin=8pt]
\centering\itshape
\og La méthode d'Euler est-elle un cas isolé,
ou bien l'exemple d'une idée plus générale\,?\fg{}
\end{mdframed}
\end{center}

La réponse que l'on va construire aujourd'hui :
l'analyse numérique repose souvent sur la même stratégie
--- remplacer des objets continus par des objets discrets,
transformer une équation en relations algébriques,
et obtenir un problème calculable.

% ====================================================================
\section*{2. Des taux de variation aux différences finies}
% ====================================================================

Au CM1, pour construire Euler, on a remplacé
la dérivée $\frac{dT}{dt}$ par un quotient :
$\frac{T_{n+1}-T_n}{\Delta t}$.
Cette idée porte un nom : les \textbf{différences finies}.
Elle s'étend à toute dérivée, en espace comme en temps.

\medskip

\textbf{Rappel.} La dérivée de $u$ en $x$ est définie comme :
\[
  u'(x) = \lim_{h \to 0} \frac{u(x+h) - u(x)}{h}.
\]

Si l'on supprime le passage à la limite
et que l'on garde $h > 0$ fixé,
on obtient une \textbf{approximation} de la dérivée.

\medskip

\textbf{Question.}
Supposons que $u$ n'est connue qu'en des points
$x_{i-1}$, $x_i$, $x_{i+1}$ régulièrement espacés de $h$.
Comment approximer $u'(x_i)$ ?

\atrou[4cm]{%
Trois approximations classiques (on note $u_j = u(x_j)$) :
\begin{itemize}
  \item \textbf{Différence progressive} (avant) :
    $\displaystyle u'(x_i) \approx \frac{u_{i+1} - u_i}{h}$
  \item \textbf{Différence rétrograde} (arrière) :
    $\displaystyle u'(x_i) \approx \frac{u_i - u_{i-1}}{h}$
  \item \textbf{Différence centrée} :
    $\displaystyle u'(x_i) \approx \frac{u_{i+1} - u_{i-1}}{2h}$
\end{itemize}
}

\medskip

\textbf{Et la dérivée seconde ?}
La dérivée seconde mesure comment la pente elle-même varie.
On sait que :
\[
  u''(x) = \lim_{h\to 0}
  \frac{u(x+h) - 2\,u(x) + u(x-h)}{h^2}.
\]

\textbf{Question.}
Comment approximer $u''(x_i)$
à partir de $u_{i-1}$, $u_i$ et $u_{i+1}$ ?

\atrou[3cm]{%
\[
  \boxed{u''(x_i) \approx
  \frac{u_{i+1} - 2\,u_i + u_{i-1}}{h^2}.}
\]
Cette formule fait intervenir trois points voisins :
c'est un \textbf{stencil à trois points}.
}

\medskip

\textbf{Interprétation graphique}
--- la dérivée seconde mesure la courbure de $u$ :

\esquisse[6cm]

% ====================================================================
\section*{3. Exemple --- De la physique au système linéaire}
% ====================================================================

On va maintenant utiliser les différences finies
pour résoudre un problème d'un type nouveau :
un \textbf{problème aux limites},
où l'inconnue est un profil spatial
et non une trajectoire temporelle.

\medskip

\begin{center}
\begin{mdframed}[backgroundcolor=ThemeLight, linecolor=Theme,
  linewidth=1.5pt, innertopmargin=10pt, innerbottommargin=10pt]
\centering\large\itshape
\og Une barre métallique, chauffée uniformément,\\
est maintenue à $0\;\degC$ à ses deux extrémités.\\
Quel est le profil de température à l'équilibre\,?\fg{}
\end{mdframed}
\end{center}

\textbf{Questions préliminaires.}
\begin{itemize}[leftmargin=*]
  \item À quoi ressemble le profil de température ?
        Faites un croquis.
  \item Si l'on isole une petite tranche de barre,
        quels échanges thermiques y a-t-il ?
\end{itemize}

\atrou[4cm]{%
\begin{itemize}[leftmargin=*]
  \item Le profil est nul aux deux bouts et maximal au milieu
        (la chaleur générée au centre doit
        parcourir la plus grande distance pour atteindre les bords).
  \item Trois contributions : le \textbf{flux de conduction
        entrant} par la gauche, le \textbf{flux sortant}
        par la droite, et la \textbf{chaleur générée}
        dans la tranche. À l'équilibre, le bilan est nul.
\end{itemize}
}

\textbf{Bilan thermique sur une tranche $[x,\,x+dx]$.}
La loi de Fourier donne le flux de chaleur :
$\varphi = -\lambda\,\dfrac{du}{dx}$.

\smallskip

Le bilan d'énergie à l'équilibre :
\[
  \underbrace{\varphi(x)\,A}_{\text{flux entrant}}
  - \underbrace{\varphi(x+dx)\,A}_{\text{flux sortant}}
  + \underbrace{q\,A\,dx}_{\text{chaleur générée}}
  = 0.
\]

En divisant par $A\,dx$ et en passant à la limite :

\atrou[2.5cm]{%
\[
  -\frac{d\varphi}{dx} + q = 0
  \qquad\Longrightarrow\qquad
  \boxed{-\lambda\,\frac{d^2 u}{dx^2} = q.}
\]
}

\textbf{Forme simplifiée.}
Après adimensionnement ($L = 1$, division par $\lambda$) :
\[
  -u''(x) = 1, \qquad x \in {]0,\,1[},
  \qquad u(0) = 0,\quad u(1) = 0.
\]

$u(x)$ représente la température en un point $x$ de la barre.

\medskip

\textbf{Étape 1 --- Maillage.}
On découpe l'intervalle $[0,1]$ en $N$ sous-intervalles
de même longueur $h = \frac{1}{N}$.
Les \textbf{nœuds} sont $x_i = i\,h$ pour $i = 0, 1, \ldots, N$.
On cherche des valeurs approchées $u_i \approx u(x_i)$.

\esquisse[4cm]

\textbf{Questions.}
\begin{itemize}[leftmargin=*]
  \item Parmi $u_0, u_1, \ldots, u_N$,
        lesquels sont déjà connus ?
  \item Combien d'inconnues reste-t-il ?
  \item Combien d'équations faut-il ?
\end{itemize}

\atrou[2.5cm]{%
$u_0 = 0$ et $u_N = 0$ sont imposés
par les conditions aux limites.
Il reste $N-1$ inconnues :
$u_1, u_2, \ldots, u_{N-1}$.
Il faut donc $N-1$ équations.
}

\medskip

\textbf{Étape 2 --- Écriture des équations.}
En chaque nœud intérieur $x_i$ ($i = 1, \ldots, N-1$),
on remplace $-u''(x_i)$ par la formule
de différences finies :

\atrou[3cm]{%
\[
  -\frac{u_{i+1} - 2\,u_i + u_{i-1}}{h^2} = 1,
  \qquad i = 1, \ldots, N-1.
\]
En multipliant par $h^2$ :
\[
  -u_{i-1} + 2\,u_i - u_{i+1} = h^2.
\]
}

\medskip

\textbf{Étape 3 --- Construction du système.}
Prenons $N = 5$
(donc $h = 0{,}2$, $h^2 = 0{,}04$
et $4$ inconnues : $u_1, u_2, u_3, u_4$).

\smallskip

\textbf{Activité :} écrire les $4$ équations
en utilisant $u_0 = u_5 = 0$.

\atrou[7cm]{%
\begin{align*}
  i=1 :\quad 2u_1 - u_2 &= 0{,}04
  && \text{(car $u_0 = 0$)} \\
  i=2 :\quad -u_1 + 2u_2 - u_3 &= 0{,}04 \\
  i=3 :\quad -u_2 + 2u_3 - u_4 &= 0{,}04 \\
  i=4 :\quad -u_3 + 2u_4 &= 0{,}04
  && \text{(car $u_5 = 0$)}
\end{align*}

\medskip
Sous forme matricielle : $A\,\mathbf{u} = \mathbf{b}$, soit :
\[
  \underbrace{\begin{pmatrix}
    2 & -1 & 0 & 0 \\
    -1 & 2 & -1 & 0 \\
    0 & -1 & 2 & -1 \\
    0 & 0 & -1 & 2
  \end{pmatrix}}_{A}
  \begin{pmatrix} u_1 \\ u_2 \\ u_3 \\ u_4 \end{pmatrix}
  =
  \begin{pmatrix} 0{,}04 \\ 0{,}04 \\ 0{,}04 \\ 0{,}04 \end{pmatrix}.
\]
}

\medskip

\textbf{Questions.}
\begin{itemize}[leftmargin=*]
  \item Quelle est la structure de la matrice $A$ ?
  \item Pourquoi n'y a-t-il que trois diagonales non nulles ?
  \item Si l'on prenait $N = 100$,
        quelle serait la taille du système ?
\end{itemize}

\atrou[3.6cm]{%
\begin{itemize}
  \item $A$ est \textbf{tridiagonale} :
        seules la diagonale principale
        et les deux diagonales adjacentes sont non nulles.
  \item Cela vient du stencil à trois points :
        chaque équation en $x_i$ ne fait intervenir
        que $u_{i-1}$, $u_i$ et $u_{i+1}$.
  \item Avec $N = 100$, le système serait
        de taille $99 \times 99$,
        mais la matrice resterait tridiagonale :
        très peu de coefficients non nuls.
\end{itemize}
}

\medskip

\textbf{Vérification.}
La solution exacte du problème continu est
$u(x) = \frac{x\,(1-x)}{2}$.

\begin{center}
\renewcommand{\arraystretch}{1.6}
\begin{tabular}{|c|c|c|c|}
  \hline
  $i$ & $x_i$ & $u_i$ (numérique) & $u(x_i) = \frac{x_i(1-x_i)}{2}$ \\
  \hline
  1 & $0{,}2$ & \val{$0{,}08$} & $0{,}08$ \\
  \hline
  2 & $0{,}4$ & \val{$0{,}12$} & $0{,}12$ \\
  \hline
  3 & $0{,}6$ & \val{$0{,}12$} & $0{,}12$ \\
  \hline
  4 & $0{,}8$ & \val{$0{,}08$} & $0{,}08$ \\
  \hline
\end{tabular}
\end{center}

\smallskip

\textit{Ici, les valeurs coïncident car la solution exacte
est un polynôme de degré~$2$
et le schéma est exact pour de tels polynômes.
Pour une source $f(x)$ plus complexe,
on observerait un écart qui diminue quand $N$ augmente.}

% ====================================================================
\section*{4. Ouverture --- Diffusion de chaleur dans une paroi}
% ====================================================================

Les mêmes idées de discrétisation s'appliquent
à des problèmes plus riches.
Reprenons un exemple de transfert thermique,
cette fois avec une dimension d'espace
\emph{et} une dimension de temps.

\medskip

La température $T(x,t)$ dans une paroi d'épaisseur $L$
vérifie l'\textbf{équation de la chaleur} :
\[
  \frac{\partial T}{\partial t}
  = \alpha\,\frac{\partial^2 T}{\partial x^2},
\]
où $\alpha > 0$ est la diffusivité thermique du matériau.

\medskip

\textbf{Lien avec ce qu'on vient de voir.}
\begin{itemize}[leftmargin=*]
  \item En \textbf{régime permanent}
        ($\frac{\partial T}{\partial t} = 0$),
        on retombe sur $T''(x) = 0$
        (ou $-T'' = f$ avec source de chaleur) :
        c'est le problème aux limites de la section~3.
  \item Si l'évolution temporelle intervient,
        on peut d'abord discrétiser en espace
        (différences finies sur $x$),
        puis avancer en temps avec Euler (CM1).
\end{itemize}

\atrou[5cm]{%
Discrétisation en espace de
$\frac{\partial^2 T}{\partial x^2}$
au nœud $x_i$ :
\[
  \frac{\partial T_i}{\partial t}
  \approx \alpha\;\frac{T_{i+1} - 2\,T_i + T_{i-1}}{h^2}.
\]

On obtient un système d'EDO couplées.
Un pas d'Euler en temps donne :
\[
  \boxed{T_i^{n+1} = T_i^n
  + \frac{\alpha\,\Delta t}{h^2}
    \bigl(T_{i+1}^n - 2\,T_i^n + T_{i-1}^n\bigr).}
\]

La même logique --- discrétisation puis calcul pas à pas ---
permet de simuler l'évolution de la température.
}

\medskip

\textbf{Question.}
À quoi reconnaît-on qu'une solution numérique
est absurde dans un contexte physique ?

\atrou[4cm]{%
Signaux d'alerte :
\begin{itemize}
  \item \textbf{Oscillations non physiques} :
        la température alterne entre chaud et froid
        d'un nœud à l'autre.
  \item \textbf{Explosion des valeurs} :
        la température diverge vers l'infini.
  \item \textbf{Incohérence avec les bords} :
        la solution ne respecte pas
        les températures imposées aux extrémités.
  \item \textbf{Comportement non compatible
        avec le phénomène} :
        la chaleur semble se propager dans le mauvais sens.
\end{itemize}
}

% ====================================================================
\section*{5. Stabilité et convergence --- Premières intuitions}
% ====================================================================

On a vu qu'en discrétisant un problème continu,
on obtient une solution approchée qui dépend
du pas de maillage.
Comment juger si cette approximation est fiable ?

\medskip

\textbf{Convergence.}

\begin{center}
\begin{mdframed}[backgroundcolor=ThemeLight, linecolor=Theme,
  linewidth=1pt, innertopmargin=8pt, innerbottommargin=8pt]
\centering\itshape
Quand le pas tend vers zéro,
la solution numérique se rapproche-t-elle
de la solution exacte\,?
\end{mdframed}
\end{center}

\atrou[3.6cm]{%
Si la méthode est \textbf{convergente},
l'écart entre solution numérique et solution exacte
tend vers zéro quand le pas diminue.

\smallskip

Exemple (CM1) : pour le café,
Euler avec $\Delta t = 5$ donne une erreur de $8{,}3\;\degC$,
avec $\Delta t = 2$ l'erreur tombe à $2{,}8\;\degC$.
En diminuant le pas, on se rapproche de la solution exacte.
}

\medskip

\textbf{Stabilité.}

\begin{center}
\begin{mdframed}[backgroundcolor=ThemeLight, linecolor=Theme,
  linewidth=1pt, innertopmargin=8pt, innerbottommargin=8pt]
\centering\itshape
Les petites erreurs de calcul restent-elles
sous contrôle, ou explosent-elles\,?
\end{mdframed}
\end{center}

\atrou[4cm]{%
Une méthode est \textbf{stable} si les erreurs
(arrondis, perturbations) ne sont pas amplifiées
de façon catastrophique au fil des itérations.

\smallskip

Exemple (CM1) : Euler sur le café
avec $\Delta t = 15$ donne $T_1 = -15\;\degC$.
L'erreur a explosé en une seule étape :
le schéma est \textbf{instable} pour ce pas.

\smallskip

Pour l'équation de la chaleur,
la condition de stabilité s'écrit :
\[
  \frac{\alpha\,\Delta t}{h^2} \leq \frac{1}{2}.
\]
Un pas de temps trop grand par rapport au maillage spatial
provoque des oscillations et la divergence.
}

\medskip

\textbf{Compromis précision / coût.}

\atrou[3.6cm]{%
Raffiner le maillage améliore l'approximation,
mais augmente le nombre d'inconnues
et le temps de calcul.
\begin{itemize}
  \item En 1D avec $N$ sous-intervalles :
        système de taille $N-1$.
  \item En 2D sur un maillage $N \times N$ :
        de l'ordre de $N^2$ inconnues.
  \item Le choix du pas est un \textbf{compromis
        entre précision et coût}.
\end{itemize}
}

\medskip

\textbf{Discussion.}
Deux maillages différents pour le même problème :

\begin{center}
\renewcommand{\arraystretch}{1.4}
\begin{tabular}{cccl}
  \hline
  Maillage & Inconnues & Coût & Précision \\
  \hline
  $N = 10$ & $9$ & \val{faible} & \val{modérée} \\
  $N = 1000$ & $999$ & \val{élevé} & \val{excellente} \\
  \hline
\end{tabular}
\end{center}

En ingénierie, quel compromis vous semblerait raisonnable ?

\atrou[3cm]{%
La réponse dépend du contexte :
\begin{itemize}
  \item \textbf{Conception préliminaire} :
        un calcul rapide avec un maillage grossier suffit
        pour explorer des variantes.
  \item \textbf{Dimensionnement final} :
        un calcul plus fin est nécessaire
        pour garantir la sécurité.
  \item Dans tous les cas, il est essentiel de
        \textbf{vérifier la plausibilité physique
        du résultat}.
\end{itemize}
}

% ====================================================================
\section*{6. Bilan}
% ====================================================================

Quatre idées à retenir :

\atrou[7cm]{%
\begin{enumerate}
  \item L'idée de \textbf{différences finies}
        est une généralisation de ce qu'on a fait avec Euler :
        on remplace des dérivées par des quotients
        de différences.
  \item La discrétisation d'un problème aux limites
        conduit naturellement à un \textbf{système linéaire}
        $A\,\mathbf{u} = \mathbf{b}$.
  \item La matrice obtenue a une structure particulière
        (\textbf{tridiagonale} en 1D)
        qui reflète le caractère local du stencil.
  \item La fiabilité d'un calcul numérique repose sur
        trois propriétés :
        \textbf{convergence}
        (le résultat s'améliore quand on raffine),
        \textbf{stabilité}
        (les erreurs ne sont pas amplifiées),
        et un \textbf{compromis entre précision et coût}.
\end{enumerate}
}

\end{document}
